%% ============================================================
%%  main.tex — IEEE Conference Paper
%%  "Uncertainty- and Impact-Aware Human-in-the-Loop Gating
%%   for Multi-Agent Software Engineering"
%%
%%  Build:  pdflatex main && bibtex main && pdflatex main && pdflatex main
%%
%%  Anonymous switch: change \anonymousfalse -> \anonymoustrue
%%  before blind submission.
%% ============================================================
\documentclass[conference]{IEEEtran}
\IEEEoverridecommandlockouts

%% -- Packages --------------------------------------------------
\usepackage{cite}
\usepackage{amsmath,amssymb,amsfonts}
\usepackage{algorithmic}
\usepackage{graphicx}
\usepackage{textcomp}
\usepackage{xcolor}
\usepackage{listings}
\usepackage{booktabs}
\usepackage{tabularx}
\usepackage{array}
\usepackage{multirow}
\usepackage{url}
\usepackage{balance}
\usepackage[colorinlistoftodos,prependcaption]{todonotes}

%% -- Anonymous switch ------------------------------------------
\newif\ifanonymous
\anonymousfalse   %% Set \anonymoustrue for blind review

%% -- Notation macros -------------------------------------------
\newcommand{\Usem}{U_{\mathrm{sem}}}
\newcommand{\Uent}{U_{\mathrm{ent}}}
\newcommand{\Uref}{U_{\mathrm{ref}}}
\newcommand{\Tbase}{T_{\mathrm{base}}}
\newcommand{\Treview}{T_{\mathrm{review}}}
\newcommand{\Tnorm}{T_{\mathrm{norm}}}
\newcommand{\Tmin}{T_{\mathrm{min}}}
\newcommand{\Tmax}{T_{\mathrm{max}}}

%% -- Listings style -------------------------------------------
\lstdefinestyle{jsonStyle}{
  basicstyle=\scriptsize\ttfamily,
  frame=single,
  breaklines=true,
  columns=fullflexible,
  keepspaces=true,
  showstringspaces=false,
}
\lstset{style=jsonStyle}

%% -- BibTeX ----------------------------------------------------
\def\BibTeX{{\rm B\kern-.05em{\sc i\kern-.025em b}\kern-.08em
    T\kern-.1667em\lower.7ex\hbox{E}\kern-.125emX}}

%% ============================================================
\begin{document}

%% -- Title ----------------------------------------------------
\title{Uncertainty- and Impact-Aware Human-in-the-Loop Gating\\
       for Multi-Agent Software Engineering}

%% -- Authors --------------------------------------------------
\ifanonymous
  \author{\IEEEauthorblockN{Anonymous Authors}
  \IEEEauthorblockA{Anonymous Institution}}
\else
\author{
\begin{minipage}{\textwidth}
\centering
\begin{tabular}{ccc}
\textbf{Yukthesh Reddy Konda} &
\textbf{Jayadeep Kundrapu} &
\textbf{Navneeth Kundrapu} \\
\textit{Dept.~of CSE} &
\textit{Dept.~of CSE} &
\textit{Dept.~of ECE} \\
\textit{IIIT Naya Raipur, India} &
\textit{IIIT Naya Raipur, India} &
\textit{IIIT Naya Raipur, India} \\
konda23100@iiitnr.edu.in &
kundrapu23100@iiitnr.edu.in &
kundrapu23101@iiitnr.edu.in \\
\end{tabular}
\vspace{0.5em}\\
{\large \textbf{Dr.~Mallikharjun Rao K}} \\
\textit{Data Science and Artificial Intelligence,
        IIIT Naya Raipur, India} \\
mallikharjuna@iiitnr.edu.in \quad
\url{https://orcid.org/0000-0002-6096-3963}
\end{minipage}
}
\fi

\maketitle

%% -- Abstract -------------------------------------------------
\begin{abstract}
Autonomous multi-agent systems that orchestrate Large Language Models (LLMs)
across a software-development life cycle are increasingly capable but remain
susceptible to hallucinated architectural decisions whose errors compound
through downstream agents.
We describe a Human-in-the-Loop (HITL) gating mechanism---the
\emph{Rejector Gate}---that intercepts each agent output and decides whether
to route execution to autonomous continuation or to request human review.
The gate computes a composite uncertainty score $U$ from three signals:
lexical dispersion across $N$ response variants, normalized decision entropy
over structured choices, and an LLM-based reflection score.
$U$ is then scaled by a decision gravity factor $G$, derived from impact,
irreversibility, and propagation of the decision.
Intervention is triggered when $S = U \times G$ exceeds an adaptive
threshold.
To set that threshold, we formulate the gating problem as an episodic Markov
Decision Process and train a Deep Q-Network (DQN) using a token-economics
reward model that analytically derives a large asymmetry between
false-negative and false-positive penalties (ranging from $3.6\times$ at
minimum gravity to $79\times$ at maximum).
A secondary finite-difference optimizer tunes the continuous weighting
coefficients on a slower timescale.
We evaluate the system on a simulation environment comprising 20 curated and
up to 50 procedurally generated architectural decision scenarios per episode.
Aggregate empirical results on real projects and comparisons against published
baselines are deferred to future work and listed in the accompanying
\texttt{TODO.md}.
\end{abstract}

\begin{IEEEkeywords}
Multi-Agent Systems, Human-in-the-Loop, Reinforcement Learning,
Large Language Models, Software Engineering, Deep Q-Network,
Uncertainty Estimation, Decision Gravity
\end{IEEEkeywords}

%% ============================================================
\section{Introduction}\label{sec:intro}
%% ============================================================

Large Language Models (LLMs) have become capable collaborators in
software-engineering tasks~\cite{brown2020gpt3,bubeck2023gpt4,hou2024llmse}.
Frameworks such as MetaGPT~\cite{hong2024metagpt} and
ChatDev~\cite{qian2024chatdev} assign specialized \emph{personas}---product
manager, architect, coder, tester---to LLMs and demonstrate that such
role-based decomposition improves code quality over monolithic prompting.
Yet purely autonomous pipelines share a structural vulnerability: when an
early-stage agent produces an incorrect architectural decision, every
downstream agent that depends on that decision inherits and amplifies the
error.
Mitigating this risk by routing every agent output to a human reviewer
would eliminate the efficiency gains that motivate
automation~\cite{horvitz1999principles,amershi2019guidelines}.

The key challenge is therefore \emph{when} to request human review.
Existing approaches rely either on static confidence thresholds that do not
adapt to operational context, or on RLHF-style training that aligns model
weights offline but does not gate runtime
decisions~\cite{christiano2017rlhf,ouyang2022instructgpt}.
Selective-prediction and learning-to-defer
methods~\cite{geifman2017selective,madras2018predict} offer formal frameworks
for abstaining, but have not been applied to multi-agent software pipelines
where the \emph{consequences} of an error---not just its probability---
determine the urgency of intervention.

\textbf{Contributions.} This paper makes three contributions:
\begin{enumerate}
  \item \textbf{The Rejector Gate:} a gating mechanism combining three
        complementary uncertainty signals (lexical dispersion, decision
        entropy, and LLM reflection) weighted by a decision gravity score
        to produce a composite gate score $S = U \times G$ that determines
        routing.
  \item \textbf{Token-economics reward model:} a reward function derived
        analytically from estimated token expenditures that produces an
        asymmetric FN/FP penalty ratio of approximately $21\times$ at medium
        gravity and $79\times$ at maximum gravity, without hand-tuned
        constants.
  \item \textbf{Dual-layer optimization:} a DQN that learns a binary
        intervention policy from simulated decisions, combined with a
        finite-difference parameter optimizer that tunes the uncertainty
        weighting coefficients on a slower timescale.
\end{enumerate}

\textbf{Paper outline.}
Section~\ref{sec:related} reviews related work.
Section~\ref{sec:arch} describes the system architecture.
Section~\ref{sec:rationale} explains the design rationale.
Section~\ref{sec:math} presents the mathematical formulation.
Section~\ref{sec:reward} derives the reward model.
Section~\ref{sec:rl} describes the learning framework.
Section~\ref{sec:eval} defines the experimental setup.
Section~\ref{sec:results} reports results.
Section~\ref{sec:threats} discusses threats to validity.
Section~\ref{sec:conclusion} concludes.

%% ============================================================
\section{Related Work}\label{sec:related}
%% ============================================================

\subsection{LLMs for Software Engineering}
Surveys such as Hou et al.~\cite{hou2024llmse} document rapid growth in
LLM-assisted programming, code review, and test generation.
SWE-agent~\cite{yang2024sweagent} demonstrates end-to-end issue resolution
in real repositories.
Our work differs in that we do not replace the developer but decide
\emph{dynamically} when their involvement is most valuable.

\subsection{Multi-Agent Frameworks}
MetaGPT~\cite{hong2024metagpt} encodes software engineering workflows as
structured LLM conversations with role-based personas.
ChatDev~\cite{qian2024chatdev} uses a waterfall-style agent chain with
chat-based inter-agent debate.
AutoGen~\cite{wu2024autogen} provides a general conversational multi-agent
framework.
All three operate with either no gating or static heuristics.
Our work adds an adaptive gate that interrupts the chain only when both
uncertainty and decision gravity are high.

\subsection{LLM Uncertainty and Hallucination Detection}
SelfCheckGPT~\cite{manakul2023selfcheckgpt} detects hallucinations by
sampling multiple responses and measuring their
consistency---the same principle underlies our $\Usem$.
Kuhn et al.~\cite{kuhn2023semantic} introduce semantic entropy, clustering
semantically equivalent outputs before computing entropy.
Kadavath et al.~\cite{kadavath2022know} show that LLMs can be calibrated
to predict their own correctness.
Our $\Uref$ adopts this idea by prompting the LLM to review its own
variants, while acknowledging the same-model-family bias this creates.

\subsection{Selective Prediction and Learning to Defer}
Geifman and El-Yaniv~\cite{geifman2017selective} formulate selective
prediction as minimizing risk subject to a coverage constraint.
Madras et al.~\cite{madras2018predict} extend this to a human-AI team where
deferral decisions consider both accuracy and fairness.
Our work is in this spirit but grounds the deferral cost in token economics
rather than a fixed abstention penalty.

\subsection{Human-AI Teaming and Mixed-Initiative Interaction}
Horvitz~\cite{horvitz1999principles} articulates principles for
mixed-initiative systems, arguing that control allocation should be
sensitive to task context.
Amershi et al.~\cite{amershi2019guidelines} provide empirical guidelines
for human-AI interaction design.
Our gate operationalises context-sensitivity through the gravity score $G$.

\subsection{RLHF and Human-in-the-Loop RL}
Christiano et al.~\cite{christiano2017rlhf} use human preferences to
shape a reward model for alignment.
Ouyang et al.~\cite{ouyang2022instructgpt} scale this to instruction-following
via PPO.
Both approaches align model weights offline; they do not gate runtime
decisions.
Our DQN learns a \emph{runtime} binary policy that decides per-output
whether to invoke the human, complementing rather than replacing offline
alignment.
Active learning~\cite{settles2012active} also queries an oracle selectively,
but typically assumes a static pool rather than a sequential pipeline.

\subsection{Comparison Summary}

Table~\ref{tab:related} summarises how our system relates to representative
prior work.

\begin{table}[!t]
  \caption{Comparison of multi-agent and gating approaches}
  \label{tab:related}
  \renewcommand{\arraystretch}{1.2}
  \resizebox{\columnwidth}{!}{%
  \begin{tabular}{@{}lp{1.5cm}p{1.5cm}p{0.9cm}p{0.9cm}p{1.0cm}@{}}
    \toprule
    \textbf{System} &
    \textbf{Human Involvement} &
    \textbf{Uncertainty} &
    \textbf{Adaptive} &
    \textbf{Cost Model} &
    \textbf{Eval.~Scale} \\
    \midrule
    MetaGPT~\cite{hong2024metagpt}        & None          & None             & No  & No         & Projects  \\
    ChatDev~\cite{qian2024chatdev}        & Static review & None             & No  & No         & Projects  \\
    AutoGen~\cite{wu2024autogen}          & Optional      & None             & No  & No         & Tasks     \\
    Madras et al.~\cite{madras2018predict}& On defer      & Confidence       & Yes & Coverage   & Benchmarks\\
    \textbf{Ours}                         & Gate-triggered& $\Usem,\Uent,\Uref$ & Yes & Token econ.& Simulation\\
    \bottomrule
  \end{tabular}}
\end{table}

%% ============================================================
\section{System Architecture}\label{sec:arch}
%% ============================================================

\subsection{Target System Class}

The framework targets CRUD web applications with a REST backend:
projects where the agent chain produces API endpoint definitions,
relational or document data models, and service-layer code.
This scope constrains what structured decisions can be extracted for entropy
computation and what constitutes a recoverable versus catastrophic
architectural error.

\subsection{Agent Pipeline}

As illustrated in Fig.~\ref{fig:architecture}, requirements flow through
four specialized agents:

\begin{enumerate}
  \item \textbf{System Architect:} ingests raw requirements and emits a
        structured JSON plan containing \texttt{projectName},
        \texttt{techStack}, a list of \texttt{components} with dependency
        edges, \texttt{apiEndpoints} (method, path, request/response schema),
        \texttt{dataModels} with typed fields, and---critically---a list of
        \texttt{architecturalDecisions}, each recording the decision
        statement, the options considered, the chosen option, and the
        agent's reasoning.
        This field is the primary input for $\Uent$.
  \item \textbf{Orchestrator:} decomposes the architectural plan into
        concrete, decoupled tasks assigned to domain-specific agents.
  \item \textbf{Frontend Agent:} generates UI code, routing, and
        schema-validation logic.
  \item \textbf{Backend Agent:} generates data-access layer code,
        API route handlers, and authentication logic.
\end{enumerate}

No ``Integrator Agent'' is present in the current implementation.

\begin{figure}[!h]
  \centering
  \includegraphics[width=\columnwidth]{figures/architecture.pdf}
  \caption{System Architecture. Requirements flow through the System
           Architect, Orchestrator, and Implementation agents.
           The Rejector Gate intercepts each agent's output before
           it is passed downstream.}
  \label{fig:architecture}
  \todo[color=orange!60,inline]{Regenerate as vector PDF/TikZ.
  Fix ``System Architech'' typo from original figure.
  Remove Integrator Agent (not implemented). See TODO.md item 11.}
\end{figure}

\subsection{The Rejector Gate}\label{subsec:rejector}

After each agent produces its primary output, the Rejector Gate intercepts
it before the next agent receives it (see Fig.~\ref{fig:rejector}).
The gate: (i)~generates $N{-}1$ additional response variants at temperature
$T{=}1.0$, forming a set of $N$~variants (default $N{=}3$);
(ii)~computes composite uncertainty $U$ from three signals;
(iii)~obtains decision gravity $G$ via an independent LLM assessment;
(iv)~computes gate score $S = U \times G$;
(v)~routes to \textsc{Collaborate} if $S > \tau$ (default $\tau = 0.15$)
or to \textsc{Automate} otherwise.

\begin{figure}[!h]
  \centering
  \includegraphics[width=\columnwidth]{figures/rejector_gate.pdf}
  \caption{Rejector Gate flow logic. The gate evaluates the agent's output
           against a threshold $\tau$ and routes to either Automate
           (proceed to next agent) or Collaborate (request human review).}
  \label{fig:rejector}
  \todo[color=orange!60,inline]{Regenerate as vector PDF/TikZ.
  See TODO.md item 12.}
\end{figure}

\subsection{Motivating Example}\label{subsec:example}

Consider the \emph{Monolith vs.\ Microservices} decision (impact $I{=}0.9$,
irreversibility $R{=}0.9$, propagation $P{=}1.0$; from preset scenario
\texttt{arch-pattern}).
Three variants are generated.
Variant~1 chooses Microservices citing scalability; Variant~2 chooses Monolith
citing team size; Variant~3 is split.
Pairwise cosine similarity of TF vectors averages $0.52$, giving
$\Usem{=}0.48$.
Parsing the \texttt{architecturalDecisions} field reveals two distinct choices,
yielding $\Uent \approx 1.00$ (maximum entropy for $n{=}2$ options).
The reflection LLM reports $\mathrm{agreement}{=}0.45$ with one high-severity
conflict (\textit{scaling strategy unclear}), giving $\Uref{=}0.55$.
With weights $(w_1,w_2,w_3){=}(0.40,0.35,0.25)$:
$U = 0.40{\times}0.48 + 0.35{\times}1.00 + 0.25{\times}0.55 = 0.68$.
Gravity $G = 0.9 \times 0.9 \times 1.0 = 0.81$.
Gate score $S = 0.68 \times 0.81 = 0.55 > 0.15 = \tau$:
\textsc{Collaborate}.
A human developer is alerted, selects Microservices with explicit scalability
annotations, and downstream agents propagate that choice consistently.

\subsection{System Architect Output}

Listing~\ref{lst:architect} shows an excerpt of the structured JSON output.
The \texttt{architecturalDecisions} array provides the structured signal for
$\Uent$.

\begin{lstlisting}[caption={Excerpt: System Architect JSON output},
                   label={lst:architect},
                   style=jsonStyle]
{
  "projectName": "eCommerce REST API",
  "techStack": {
    "frontend": "React",
    "backend": "NestJS",
    "database": "PostgreSQL"
  },
  "apiEndpoints": [
    { "method": "POST", "path": "/api/orders",
      "requestBody": "OrderDto",
      "responseBody": "OrderEntity" }
  ],
  "architecturalDecisions": [
    {
      "decision": "Backend framework",
      "options": ["Express", "NestJS"],
      "chosen": "NestJS",
      "reasoning": "Structured DI container reduces coupling",
      "category": "architecture"
    }
  ]
}
\end{lstlisting}

%% ============================================================
\section{Design Rationale}\label{sec:rationale}
%% ============================================================

\textbf{Why three uncertainty signals?}
Each signal captures a structurally different failure mode.
Lexical dispersion ($\Usem$) detects when the agent produces responses whose
surface forms diverge widely---a symptom of ambiguous or underdetermined
prompts.
Decision entropy ($\Uent$) captures inconsistency in the \emph{discrete}
architectural choices themselves across variants, independent of lexical
style.
The reflection score ($\Uref$) uses an LLM to identify semantic
contradictions and factual conflicts that the variance-based metrics may miss
when variants are lexically similar but logically contradictory.

\textbf{Why decision gravity?}
Uncertainty alone does not determine whether human review is warranted.
A fully reversible formatting choice may have high $\Usem$ without requiring
intervention; a schema migration with unanimous but incorrect LLM agreement
may have low $U$ but severe consequences.
Multiplying by $G$ ensures that only high-consequence uncertain decisions
trigger interruptions.

\textbf{Why RL?}
A fixed threshold $\tau$ cannot adapt to shifts in the scenario distribution
or to evolving LLM calibration.
The DQN learns which $(U, G, S)$ configurations lead to positive outcomes
across a large simulated distribution.
Without RL, the system functions with fixed $\tau{=}0.15$; the RL layer
adds adaptive calibration.
\todo[color=red!40]{Run ablation comparing fixed $\tau$ vs.~RL policy.
See TODO.md item 3.}

\textbf{Why are FP/TN human-side and TP/FN system-side?}
When the gate triggers unnecessarily (FP) or correctly withholds a safe
decision (TN), a reviewer pays a fixed cost $\Treview$~tokens regardless of
architectural significance.
When the gate correctly catches a problem (TP) or misses one (FN), the cost
scales with the complexity and reversibility of the task.
This asymmetry is encoded in the reward model (Section~\ref{sec:reward}).

%% ============================================================
\section{Mathematical Formulation}\label{sec:math}
%% ============================================================

\subsection{Lexical Dispersion (\texorpdfstring{$\Usem$}{Usem})}

Each of the $N$ variant texts is converted to a term-frequency (TF) sparse
vector $\mathbf{v}_i$ by tokenising on whitespace, discarding tokens of
length $\leq 2$, and normalising counts by total token count.
Pairwise cosine similarity is computed for all $\binom{N}{2}$ pairs:
\begin{equation}
  \label{eq:usem}
  \Usem = 1 - \frac{1}{\binom{N}{2}}
          \sum_{i < j}
          \frac{\mathbf{v}_i \cdot \mathbf{v}_j}
               {\|\mathbf{v}_i\|\,\|\mathbf{v}_j\|}
\end{equation}
For $N{=}3$ this yields three pairs.
$\Usem \in [0,1]$ by construction: identical responses give $\Usem{=}0$;
completely disjoint vocabularies give $\Usem{=}1$.

\noindent\textit{Note on simulation vs.~live pipeline:}
The RL simulation module uses 8-dimensional unit-normalised random vectors
with controlled perturbation noise as a proxy for embeddings (see
\texttt{src/simulation/scenarios.js:17--32}).
Equation~\eqref{eq:usem} describes the live gating pipeline in
\texttt{src/agents/rejector.js:66--84}.

\subsection{Decision Entropy (\texorpdfstring{$\Uent$}{Uent})}

The Rejector parses \texttt{architecturalDecisions[$k$].chosen} from each
variant.
Let $k = 1,\ldots,K$ index distinct decision points and $n_k$ denote the
number of unique choices observed across $N$ variants for point $k$.
\begin{equation}
  \label{eq:uent}
  \Uent = \frac{1}{K} \sum_{k=1}^{K}
          \frac{-\sum_{i=1}^{n_k} p_{ki} \log_2 p_{ki}}{\log_2 n_k}
\end{equation}
where $p_{ki}$ is the fraction of variants choosing option $i$ at point $k$.
When $n_k{=}1$ (unanimous), the term is zero and is excluded, avoiding
division by zero.
When no decision points can be parsed, $\Uent{=}0$.

\subsection{Reflection Score (\texorpdfstring{$\Uref$}{Uref})}

An independent LLM call (model: \texttt{gpt-4o-mini}, $T{=}0.2$) reviews all
$N$ variants using the prompt:
\begin{quote}\footnotesize
``Compare these $N$ responses from agent \textit{name}.
Return JSON: \texttt{\{agreement: float, conflicts: [...],
severity\_scores: [...]\}}.''
\end{quote}
Let $a \in [0,1]$ be the reported agreement and $\bar{s}$ the mean conflict
severity.
\begin{equation}
  \label{eq:uref}
  \Uref = \max\!\left(1 - a,\; \bar{s}\,(1-a)\right)
\end{equation}
\textit{Limitation:} the reviewing model is in the same model family as the
generation model; shared biases may cause the reviewer to overlook systematic
errors (see Section~\ref{sec:threats}).

\subsection{Composite Uncertainty}
\begin{equation}
  \label{eq:composite}
  U = w_1 \Usem + w_2 \Uent + w_3 \Uref, \qquad
  \sum_{i=1}^{3} w_i = 1
\end{equation}
Default initial weights: $(w_1, w_2, w_3) = (0.40, 0.35, 0.25)$.

\textbf{Agent-name heuristic (current implementation).}
The current implementation (\texttt{rejector.js:352--358}) applies
hard-coded per-agent-role adjustments: for the System Architect, $U$ is
shifted by $+0.3$ and clamped to $[0.45, 0.55]$; for implementation agents,
$U$ is scaled by $0.3$ and clamped to $[0.03, 0.20]$.
This embeds a prior that strategic decisions are more uncertain than
implementation decisions.
Any ``Strategic Entropy'' pattern observed across agent roles in the current
system is an artefact of this prior, not an emergent signal.
\todo[color=red!40]{Remove heuristic clamps and re-evaluate whether
$\Usem,\Uent,\Uref$ naturally stratify by agent role.
See TODO.md item 6.}

\subsection{Decision Gravity}

Gravity $G$ is assessed by an LLM call (\texttt{gpt-4o-mini}, $T{=}0.2$)
that rates three dimensions of the agent's output:

\begin{itemize}
  \item \textbf{Impact ($I$):} severity of potential failure (0 = cosmetic
        defect; 1 = data corruption or service outage).
  \item \textbf{Irreversibility ($R$):} difficulty of correcting post-deploy
        (0 = single config change; 1 = data loss or broken contracts).
  \item \textbf{Propagation ($P$):} blast radius across the codebase
        (0 = single function; 1 = entire architecture).
\end{itemize}

The current implementation computes:
\begin{equation}
  \label{eq:gravity}
  G = I \times R \times P
\end{equation}

\noindent\textbf{Known limitation:}
If any single factor is zero, $G$ collapses to zero regardless of others.
A floored geometric mean $G = (\max(I,\varepsilon) \cdot \max(R,\varepsilon)
\cdot \max(P,\varepsilon))^{1/3}$ with $\varepsilon{=}0.05$ would be more
robust~\cite{geifman2017selective}.
\todo[color=orange!60]{Update \texttt{gravityCalc.js} to use floored
geometric mean. See TODO.md item 7.}

\subsection{Gate Score and Decision Rule}
\begin{equation}
  \label{eq:gate}
  S = U \times G, \qquad
  \mathrm{action} =
  \begin{cases}
    \textsc{Collaborate} & S > \tau \\
    \textsc{Automate}    & \text{otherwise}
  \end{cases}
\end{equation}
Default $\tau = 0.15$ (from \texttt{rejector.js:322}).

%% ============================================================
\section{Token-Economics Reward Model}\label{sec:reward}
%% ============================================================

\subsection{Cost Asymmetry}

Table~\ref{tab:cost_asymmetry} distinguishes two outcome classes by which
party bears the cost and how that cost scales with gravity.

\begin{table}[!t]
  \caption{Cost asymmetry between human-effort and system-risk outcomes}
  \label{tab:cost_asymmetry}
  \renewcommand{\arraystretch}{1.2}
  \begin{tabular}{|>{\centering\bfseries\small}m{1.9cm}|
                   >{\small}m{5.6cm}|}
    \hline
    Cost class & What the cost represents \\
    \hline
    Human-effort outcomes (TN, FP) &
      Reviewer pays a fixed overhead $\Treview$ per inspection, independent
      of architectural severity.
      Excessive FPs increase reviewer fatigue without proportional benefit.
      Gravity $G$ does not scale this cost. \\
    \hline
    System-risk outcomes (TP, FN) &
      Cost scales with task complexity $\Tbase$ and rework multiplier
      $\kappa(R)$.
      Missing a high-gravity error (FN) is far more costly than missing a
      low-gravity one; catching it (TP) correspondingly averts more waste.
      Gravity $G$ determines $\Tbase$ and $\kappa$. \\
    \hline
  \end{tabular}
\end{table}

\subsection{Parameter Definitions}

All parameters in Table~\ref{tab:params} are taken directly from
\texttt{src/rl/tokenCostModel.js:32--36}.

\begin{table}[!t]
  \caption{Token-cost model parameters (source: \texttt{tokenCostModel.js})}
  \label{tab:params}
  \renewcommand{\arraystretch}{1.1}
  \begin{tabular}{@{}cllp{3.2cm}@{}}
    \toprule
    Symbol & Value & Unit & Rationale \\
    \midrule
    $\Tmin$    & 500      & tokens & Minimal task (config flag) \\
    $\Tmax$    & 10{,}000 & tokens & Full architectural refactor \\
    $\Treview$ & 800      & tokens & Human read + correction overhead \\
    $\alpha$   & 2.5      & ---    & Max rework ceiling multiplier \\
    $\eta$     & 0.6      & ---    & TN efficiency discount \\
    $\Tnorm$   & 45{,}000 & tokens & $\Tmax(2{+}\alpha)$ \\
    \bottomrule
  \end{tabular}
\end{table}

\subsection{Reward Formulation}

Let $\Phi(I,P) = (I+P)/2$.
The base token cost is:
\begin{equation}
  \label{eq:tbase}
  \Tbase(I,P) = \Tmin + (\Tmax - \Tmin)\,\Phi(I,P)
\end{equation}
The rework multiplier accounts for fixing a flawed implementation:
\begin{equation}
  \label{eq:kappa}
  \kappa(R) = 1 + \alpha \cdot R
\end{equation}
The four reward values, normalised by $\Tnorm$, are given in
Table~\ref{tab:rewards}.

\begin{table}[!t]
  \caption{Token-cost based reward formulation for each gate outcome}
  \label{tab:rewards}
  \renewcommand{\arraystretch}{1.3}
  \begin{tabularx}{\columnwidth}{@{}lXp{2.2cm}@{}}
    \toprule
    \textbf{Outcome} & \textbf{Reward} & \textbf{Interpretation} \\
    \midrule
    True Positive (TP) &
      $\displaystyle+\frac{\Tbase}{\Tnorm}$ &
      Tokens saved by catching the error \\[6pt]
    False Negative (FN) &
      $\displaystyle-\frac{\Tbase(1+\kappa)}{\Tnorm}$ &
      Wasted work plus rework cost \\[6pt]
    True Negative (TN) &
      $\displaystyle+\frac{\Tbase\cdot\eta}{\Tnorm}$ &
      Efficient autonomous execution \\[6pt]
    False Positive (FP) &
      $\displaystyle-\frac{\Treview}{\Tnorm}$ &
      Unnecessary review overhead \\
    \bottomrule
  \end{tabularx}
\end{table}

\subsection{FN/FP Penalty Asymmetry}\label{subsec:asymmetry}

Using $\Treview{=}800$, $\Tnorm{=}45{,}000$, and
$|R_{\mathrm{FP}}| = 800/45{,}000 \approx 0.0178$:

\begin{itemize}
  \item \textbf{Low gravity} ($I{=}R{=}P{=}0.1$):
        $\Tbase{=}1{,}450$, $\kappa{=}1.25$,
        ratio $\approx 3.6\times$
  \item \textbf{Medium gravity} ($I{=}R{=}P{=}0.5$):
        $\Tbase{=}5{,}250$, $\kappa{=}2.25$,
        ratio $\approx 21.3\times$
  \item \textbf{High gravity} ($I{=}R{=}P{=}1.0$):
        $\Tbase{=}10{,}000$, $\kappa{=}3.5$,
        ratio $\approx 79\times$
\end{itemize}

This graduated asymmetry encodes a safety-first preference: missing a
catastrophic error is penalised far more than an unnecessary review, without
relying on hand-tuned penalty constants.

%% ============================================================
\section{Learning Framework}\label{sec:rl}
%% ============================================================

\subsection{MDP Formulation}

We model gating as an episodic MDP.
One episode consists of evaluating a shuffled batch of $M$ scenarios
sequentially.

\textbf{State space:}
\begin{equation}
  \label{eq:state}
  s_t = [\Usem,\;\Uent,\;\Uref,\;I,\;R,\;P,\;S] \in \mathbb{R}^7
\end{equation}
Each component is normalised to $[0,1]$.

\textbf{Action space:}
$a_t \in \{0 : \textsc{Automate},\; 1 : \textsc{Collaborate}\}$

\textbf{Reward:} $r_t$ is selected from Table~\ref{tab:rewards} based on
the classification of $(a_t, \mathrm{gt}_t)$.

\textbf{Discount:} $\gamma = 0.99$.
Although consecutive scenarios in a shuffled episode have weak temporal
dependence, the full MDP formulation allows the agent to internalise
long-horizon trade-offs between intervention rate and accumulated token cost.

\subsection{Layer 1: Deep Q-Network}

The intervention policy is learned by a DQN~\cite{mnih2015dqn}.
Table~\ref{tab:dqn} lists all hyperparameters, verified from
\texttt{src/rl/dqnAgent.js}.

\begin{table}[!t]
  \caption{DQN hyperparameters (source: \texttt{dqnAgent.js})}
  \label{tab:dqn}
  \renewcommand{\arraystretch}{1.1}
  \begin{tabular}{@{}ll@{}}
    \toprule
    Parameter & Value \\
    \midrule
    Discount $\gamma$               & 0.99 \\
    Learning rate                   & 0.001 \\
    Batch size                      & 64 \\
    Replay buffer capacity          & 10{,}000 \\
    Target network update period    & every 100 steps \\
    Soft-update coefficient $\tau_{\rm soft}$ & 0.01 \\
    $\varepsilon$ schedule          & $1.0 \to 0.05$, ${\times}0.995$/step \\
    Hidden layers                   & [64, 32] \\
    Activations                     & ReLU (hidden), Linear (output) \\
    Initialisation                  & Xavier: $\sqrt{2/(n_{\rm in}+n_{\rm out})}$ \\
    Optimiser                       & Adam ($\beta_1{=}0.9$, $\beta_2{=}0.999$, $\varepsilon{=}10^{-8}$) \\
    Loss                            & Huber ($\delta{=}1.0$) \\
    Training episodes               & 500 \\
    \bottomrule
  \end{tabular}
\end{table}

The Bellman target is:
$\hat{Q}(s_t,a_t) = r_t + \gamma \max_{a'} Q^{-}(s_{t+1}, a'; \theta^-)$,
where $\theta^-$ are target-network parameters.

\textbf{Relationship to threshold $\tau$:}
In the live pipeline, routing uses $S > \tau$.
In simulation, the DQN directly selects action $a_t$, learning which
$(U,G,S)$ configurations should trigger intervention.
The learned decision boundary on $S$ can be recovered by probing
$Q(s,a)$ at varying score levels (implemented in
\texttt{dqnAgent.js:163--206}).

\subsection{Layer 2: Finite-Difference Parameter Optimisation}
\label{subsec:fd}

Operating every 25 episodes, a secondary optimiser tunes the pre-softmax
logits $\boldsymbol{\theta} \in \mathbb{R}^3$ for $(w_1,w_2,w_3)$.

\textbf{Method:}
This is a \emph{finite-difference / evolution-strategy} gradient
estimator~\cite{settles2012active}, not true REINFORCE~\cite{williams1992reinforce}
(which requires a stochastic policy with log-probability gradients).
Central differences:
\begin{equation}
  \label{eq:fd}
  \frac{\partial J}{\partial \theta_i} \approx
  \frac{J(\boldsymbol{\theta} + \delta\mathbf{e}_i)
      - J(\boldsymbol{\theta} - \delta\mathbf{e}_i)}{2\delta}
\end{equation}
where $J$ is the episode return evaluated by running the current DQN
greedily, perturbation $\delta{=}0.05$.
Gradients are applied with momentum ($\beta{=}0.9$, lr$=0.01$).
Weights are obtained via softmax to enforce $\sum w_i{=}1$.

\textbf{Note on $\lambda$:}
The optimiser also maintains a parameter $\lambda$ that is updated but not
consumed by the reward model in the current implementation.
\todo[color=red!40]{Either wire $\lambda$ into the reward (e.g., as a
multiplier on $\Treview$) or remove it.
See TODO.md item 8.}

%% ============================================================
\section{Experimental Setup}\label{sec:eval}
%% ============================================================

\subsection{Research Questions}

\begin{description}
  \item[RQ1] Does the RL gate reduce the false-negative rate compared to
        fixed-threshold and always-automate baselines?
  \item[RQ2] What is the intervention-rate / false-negative-rate trade-off
        curve for different gate configurations?
  \item[RQ3] Which uncertainty component contributes most to gate performance?
        (Ablation over individual signals and their combinations.)
  \item[RQ4] Does removing the agent-name heuristic clamp change performance,
        and does the strategic-entropy pattern persist?
  \item[RQ5] How does the learned decision boundary vary across decision
        categories (Architecture, Security, Infrastructure, Frontend)?
\end{description}

\subsection{Simulation Environment}

The environment (\texttt{src/rl/environment.js}) presents scenarios one at a
time to the DQN.
Each episode uses a shuffled batch of:
\begin{itemize}
  \item \textbf{20 preset scenarios:} curated architectural decisions with
        manually assigned $(I, R, P)$ and human-readable ground-truth labels
        (8 require human review; 12 do not).
        These labels serve as \emph{simulator reference labels} assigned by
        the authors; they are not independent human annotations.
  \item \textbf{50 procedurally generated scenarios:} drawn from a seeded
        pseudo-random generator (Mulberry32).
        Ground truth uses the heuristic:
        $\mathrm{needsHuman} = (G{>}0.3 \land a{<}0.7) \lor
        (G{>}0.5 \land \mathrm{rng}{>}0.5)$,
        where $a$ is LLM agreement and $G = I{\times}R{\times}P$.
\end{itemize}

\noindent\textbf{Reproducibility note:}
In the current implementation, the random batch is seeded with
\texttt{Date.now()}, making episodes irreproducible across runs.
\todo[color=red!40]{Fix PRNG seed to a constant (e.g., 42) for evaluation.
See TODO.md item 5.}

\subsection{System Under Test}

\begin{itemize}
  \item Runtime: Node.js (ESM modules, \texttt{type: "module"})
  \item LLM: \texttt{gpt-4o-mini} (default; configurable via
        \texttt{OPENAI\_MODEL} env variable; source: \texttt{llmClient.js:9})
  \item Generation temperature: $T{=}1.0$ (fixed by provider constraint;
        source: \texttt{llmClient.js:39})
  \item Reflection/gravity temperature: $T{=}0.2$
  \item $N{=}3$ variants per gate evaluation
\end{itemize}

\subsection{Baselines}

\begin{enumerate}
  \item \textbf{Always Automate:} zero interventions (lower bound on human effort)
  \item \textbf{Always Intervene:} all outputs reviewed (upper bound on safety)
  \item \textbf{Static threshold on $U$ alone} ($\tau_U$)
  \item \textbf{Static threshold on $S$} ($\tau_S{=}0.15$)
  \item \textbf{Our RL gate} (DQN + FD parameter optimizer)
  \item MetaGPT / ChatDev without gate
        \todo[color=orange!60]{Requires external setup. See TODO.md item 4.}
\end{enumerate}

\subsection{Metrics}

False-negative rate (FNR), false-positive rate (FPR), precision, recall,
F1, intervention rate, automation rate, total token cost including gate
overhead ($N$ generation calls and one reflection call per evaluation),
95\% confidence intervals, and paired Wilcoxon significance test.
\todo[color=red!40]{No result logs exist. Run experiments and populate
Section~\ref{sec:results}. See TODO.md item 2.}

%% ============================================================
\section{Results}\label{sec:results}
%% ============================================================

\subsection{Illustrative Case Study}

Table~\ref{tab:case_study} presents a single-project walkthrough of the
Rejector Gate.
This is an \emph{illustrative example}---not a systematic experiment---intended
to show the range of gate scores produced by the pipeline.
All six values are from one project run; no statistical claims are drawn
from this table alone.

\begin{table}[!t]
  \caption{Illustrative case study: single-project walkthrough
           ($\tau = 0.15$)}
  \label{tab:case_study}
  \renewcommand{\arraystretch}{1.2}
  \resizebox{\columnwidth}{!}{%
  \begin{tabular}{@{}lccclp{2.6cm}@{}}
    \toprule
    \textbf{Agent} & $U$ & $G$ & $S$ & \textbf{Action} & \textbf{Outcome} \\
    \midrule
    System Architect & 0.496 & 0.390 & 0.194 & Collaborate &
      Resolved Express vs.~NestJS conflict \\
    Orchestrator     & 0.413 & 0.459 & 0.190 & Collaborate &
      Aligned project structure and API paths \\
    Frontend (FE-1)  & 0.082 & 0.300 & 0.025 & Automate &
      Generated routing shell \\
    Frontend (FE-2)  & 0.031 & 0.280 & 0.009 & Automate &
      Generated Zod validation schemas \\
    Backend (BE-1)   & 0.180 & 0.420 & 0.075 & Automate &
      Generated Prisma/DB layer \\
    Backend (BE-2)   & 0.100 & 0.390 & 0.039 & Automate &
      Generated authentication logic \\
    \bottomrule
  \end{tabular}}
\end{table}

The pattern visible in Table~\ref{tab:case_study}---higher $U$ for strategic
agents---is consistent with the agent-name heuristic override
(Section~\ref{sec:math}).
Whether the pattern emerges without the override is deferred to RQ4.

\subsection{Aggregate Performance}
\todo[color=red!40]{Run 500-episode training, log per-episode F1, FNR, FPR,
intervention rate. Report mean $\pm$ 95\% CI for each baseline.
See TODO.md items 2 and 4.}

\subsection{Ablation Study}
\todo[color=red!40]{Train and evaluate 7 ablation variants:
each uncertainty signal alone, all three without $G$, full system,
product vs.~floored geometric mean for $G$.
See TODO.md item 3.}

\subsection{Consensus Hallucination}

A known failure mode is \emph{consensus hallucination}: all $N$ variants
converge on an incorrect or non-existent choice, yielding low $\Usem$ and
low $\Uent$.
In such cases, only $\Uref$ (the LLM reviewer) or the gravity score $G$
can prevent silent automation.
When $G$ is high, $S$ remains above $\tau$ even if $U$ is low.
However, $G$ does not itself detect the hallucination; it enforces review
when stakes are high, increasing the chance a human catches the error.
Eliminating consensus hallucination requires improving $\Uref$ or using
a cross-model reviewer.

%% ============================================================
\section{Threats to Validity}\label{sec:threats}
%% ============================================================

\textbf{Internal validity.}
Ground-truth labels for preset scenarios were assigned by the authors, not
independent annotators, and no inter-rater agreement is reported.
Ground truth for procedurally generated scenarios is a heuristic rule.
Both sources risk circular evaluation.

\textbf{External validity.}
All quantitative results derive from simulation.
The system targets CRUD REST applications; results may not generalise to
other domains.
The LLM is \texttt{gpt-4o-mini}; gravity and reflection assessments will
differ with other models.
Temperature is fixed at 1.0 by provider constraints.

\textbf{Construct validity.}
Lexical dispersion captures surface-form variance, not semantic variance
in the sense of Kuhn et al.~\cite{kuhn2023semantic}.
Decision entropy requires parseable JSON; non-compliant outputs fall back
to heuristics that may be inaccurate.
Gravity is self-assessed by the LLM producing the output, introducing
potential overconfidence or recency bias.

\textbf{Simulation bias.}
Preset scenario parameters $(I,R,P)$ reflect the authors' domain judgments.
Random parameters are sampled uniformly over $[0,1]$, which may not reflect
real project distributions.

\textbf{Same-model reflection.}
Using the same model family (\texttt{gpt-4o-mini}) to generate and to
critique variants may cause systematic blind spots where the reviewer
replicates the generator's errors.

\textbf{Perfect-human assumption.}
When intervention is triggered, the system assumes the human correction is
correct.
In practice, reviewers may introduce errors under cognitive fatigue from
excessive false positives.

%% ============================================================
\section{Conclusion}\label{sec:conclusion}
%% ============================================================

We presented a Human-in-the-Loop gating framework for multi-agent software
engineering combining three complementary uncertainty signals---lexical
dispersion, decision entropy, and LLM reflection---weighted by a decision
gravity factor to produce a composite gate score $S = U \times G$.
A Deep Q-Network, trained against a token-economics reward model with
analytically derived FN/FP asymmetry (ranging from $3.6\times$ at minimum
gravity to $79\times$ at maximum), learns a binary intervention policy over
a simulated distribution of 70 architectural decision scenarios per episode.
A finite-difference parameter optimizer tunes the uncertainty weighting
coefficients on a slower timescale.

Several important limitations remain open: the agent-name heuristic makes
the strategic-entropy pattern a tautology in the current implementation;
the gravity formula is vulnerable to zero-valued factors; and no controlled
comparison against published baselines has been conducted.
Collecting independent ground-truth labels, running systematic ablation
experiments, removing the agent-name heuristic, and validating on real
software projects constitute the principal directions for future work.

%% -- Data and Code Availability -------------------------------
\section*{Data and Code Availability}
\ifanonymous
  Code and simulation data will be made available upon acceptance.
\else
  The simulation code is available at:
  \url{https://github.com/c3-kyr/Agentic-Software-Development-with-Guardrails}
\fi

%% -- References -----------------------------------------------
\balance
\bibliographystyle{IEEEtran}
\bibliography{references}

\end{document}
